The pairing of two lifted symbols is computed from the unit parts of factorials in the prescribed digit classes. The identities below give the
precision for that pairing in the proof of
Proposition~\ref{prop:residue-plane-neighbor-test}.
\paragraph{Factorials and divided derivatives in a digit class.}
\label{par:plane-digits}
\hypertarget{source-plane-digits}{}
We record two uniform identities used in constructing the signed-chain
pairing. Write \(P=p^g\) and \(0\le\rho<P\), and use the residue
class \(\rho+P\beta\) for integers \(\beta\ge0\). Congruences in
this paragraph are coefficientwise over \(\mathbb Z_{(p)}\) and
therefore hold over the coefficient DVR as well.

For \(M\ge0\), let \(U(M)=\prod_{1\le v\le M,\,p\nmid v}v\),
with \(U(0)=1\). Removing the multiples of \(p\) from a factorial
and repeating \(g\) times gives
\begin{align*}
 (\rho+P\beta)!&=\beta!\,p^{A_\rho+\beta B_g}
       \prod_{h=0}^{g-1} U(b_h+P_h\beta),\\
 b_h&=\left\lfloor\frac\rho{p^h}\right\rfloor,\quad
 P_h=p^{g-h},\quad A_\rho=\sum_{h=0}^{g-1}\left\lfloor\frac{b_h}{p}\right\rfloor,
 \quad B_g=\frac{P-1}{p-1}.
\end{align*}
For each \(h\) put \(C_h=U(P_h)\). Since \(b_h<P_h\), splitting
into full blocks followed by the last partial block gives
\begin{align*}
 U(b_h+P_h\beta)={}&U(b_h)C_h^\beta\,
 \prod_{j=0}^{\beta-1}\prod_{\substack{1\le v\le P_h\\p\nmid v}}
                (1+P_hj/v)\,
 \prod_{\substack{1\le v\le b_h\\p\nmid v}}(1+P_h\beta/v).
\end{align*}
The two products on the right together form a unit. In a group
of coordinates with the same \(g\), the powers of \(p\) and \(C_h\)
contribute a constant depending on the total \(\sum_a\beta_a\)
and the fixed residues. The remaining unit is the product of
the displayed products of \(1\) plus a multiple of \(p\).
This gives the factorial factorization with the exact degree
of the quotient monomials as the argument of the constant.

\begin{lemma}[Factorial units and divided derivatives at fixed digits]\label{lem:plane-digit-units-and-cuts}
On a fixed digit class \(\rho+P\beta\), with \(P=p^g\),
\(g\ge0\), \(0\le\rho<P\), \(\beta\ge0\), and \(c\ge1\), either
sign of the variable unit factor in the factorial is congruent
modulo \(p^c\) to a product-binomial polynomial of total degree
at most \(2(c-1)\). A system of divided derivatives of
orders \(M_a\) admits the quotient expansion
\eqref{eq:plane-digit-derivative-expansion} with the following bound when the common radix is \(P\),
\[
 \sum_a\kappa_a\le\frac1P\sum_a M_a+c-1.
\]
 In a signed cut a divided derivative of order
\(|M|\) loses at most \(|M|\) on a vector letter and zero
on a covector letter. The factorial formula and expansions
include the radix-one case.
\end{lemma}

\begin{proof}
\proofstep{Uniform polynomial descriptions of units and inverse units.}
For \(c\ge1\), the remaining unit and its inverse both have binomial
polynomial approximations to precision \(p^c\) of total
degree at most \(2(c-1)\). To prove it, replace the leading
factor \(p\) in each \(P_hj/v\) and \(P_h\beta/v\) by a variable
\(z\). For either sign, expand the factors to degree \(c-1\)
as polynomials in \(z\) modulo \(z^c\).

In a full block at \(j\)
the coefficient of \(z^d\) is a scalar in \(\mathbb Z_{(p)}\)
times \(j^d\). Multiplying over \(j\) expresses the coefficient
of \(z^d\) as a combination of sums
\[
 \sum_{0\le j_1<\cdots<j_l<\beta}j_1^{d_1}\cdots j_l^{d_l},
 \qquad d_i\ge1,\quad \sum_i d_i=d.
\]
Induct on the number \(l\) of indices. For the induction step,
fix the last index \(j_l\). The sum over its predecessors is
an integer-valued polynomial in \(j_l\) of degree at most
\((d-d_l)+(l-1)\). Multiplication by \(j_l^{d_l}\) gives
degree at most \(d+l-1\). Summing from zero to \(\beta-1\)
raises the degree by at most one, since
\[
 \sum_{j<\beta}\binom jt=\binom{\beta}{t+1}.
\]

This proves the degree bound \(d+l\) and integrality of
the values at nonnegative integers. A partial block contributes an ordinary
polynomial of degree \(d\) in \(\beta\) at \(z^d\). The coefficient
at total order \(d\) over all blocks and coordinates is thus a
polynomial of total degree at most \(2d\), taking values in
\(\mathbb Z_{(p)}\) at nonnegative integer tuples. Write \(P(\boldsymbol\beta)\) for this coefficient polynomial and
\(\Delta_i\) for forward difference in coordinate \(i\). Its Newton
expansion is
\[
 P(\boldsymbol\beta)=\sum_{|\mathbf h|\le2d}
    (\Delta_1^{h_1}\cdots\Delta_m^{h_m}P)(0)
    \prod_{i=1}^{m}\binom{\beta_i}{h_i},
\]
where \(m\) is the number of coordinates. Each coefficient is an
integral linear combination of values of \(P\), hence lies in
\(\mathbb Z_{(p)}\). Now
put \(z=p\). This proves the approximation uniformly in the
number of coordinates and quotient exponents, also when the
fixed residues and powers \(g\) differ by group. If \(g=0\)
its product is empty and the same formulas give unit one.

\proofstep{Divided derivatives and their factorial exponents.}
For the divided derivative, write \(\partial^{[M]}X^n=\binom nM
X^{n-M}\), with the derivative zero for \(n<M\). Use
\[
 (1+T)^P=1+T^P+pE(T),\qquad \deg E\le P-1,
 \qquad E\in\mathbb Z[T].
\]
Then \(\binom{\rho+P\beta}{M}\) is the coefficient of \(T^M\) in
\((1+T)^\rho(1+T^P+pE(T))^\beta\). At order \(\ell<c\) in \(pE\),
the part with \(T^{Pj}\) from the other factor has coefficient
\[
 p^\ell[T^{M-Pj}](1+T)^\rho E(T)^\ell\,
       \binom{\ell+j}{\ell}\binom\beta{\ell+j}.
\]
This uses the binomial identity for selecting \(\ell\) marked
factors and then \(j\) others. The support of this expression
requires \(j\le M/P\) and \(M-Pj\le\rho+\ell(P-1)\).

Put \(t=\lceil(M-\rho)/P\rceil\) and \(\rho'=\rho-M+Pt\), so
\(0\le\rho'<P\). With \(\kappa=\ell+j\), the support condition
gives \(\kappa\ge t\). Therefore the corresponding divided
derivative term has the form
\begin{equation}\label{eq:plane-digit-derivative-expansion}
 \begin{aligned}
 &c_{\ell,j}X^{\rho'}
       \left(X^{\kappa-t}\partial^{[\kappa]}X^\beta\right)(X^P),\\
 &c_{\ell,j}=p^\ell[T^{M-Pj}](1+T)^\rho E(T)^\ell
                                      \binom{\ell+j}{\ell}.
 \end{aligned}
\end{equation}
All constants are integral and independent of \(\beta\).
Terms with \(\beta<\kappa\) vanish. For multiple coordinates
multiply these expansions and keep \(\sum_a\ell_a<c\), so
\[
 \sum_a\kappa_a\le\sum_a M_a/P+c-1.
\]
The discarded coefficients
are divisible by \(p^c\), also after grouping equal monomials.
This is the required derivative expansion inside the old residue
class, with the new residues \(\rho'_a\) fixed by \(M_a\).
The case \(P=1\) has \(E=0\) and gives the ordinary divided
monomial derivative exactly.

\proofstep{Effect of a derivative on a signed cut.}
In the same argument a prefix lower bound is the least Laurent
exponent for the conjugates of \(\operatorname{diag}(T,1,\ldots,1)\)
acting on the letters in the cut. For any rank-one projection
\(P_1\), write this matrix as
\[
 A_T=1+(T-1)P_1.
\]

The inverse is \(1+(T^{-1}-1)P_1\). A vector coordinate in the
cut has least order at least zero and a covector coordinate
has least order at least \(-1\). A divided derivative of
order \(|M|\) in one letter loses at most \(|M|\) in the vector
case and zero in the covector case. To see this, write the
substitution of that letter as \(x\mapsto B_Tx\). Expand the
polynomial after substitution at \(x+B_T^{-1}h\) and take the
coefficient of \(h^M\). Thus the Laurent orders of the entries satisfy
\[
 \operatorname{ord}_T\!\big((B_T^{-1})_{ij}\big)\ge\begin{cases}-1,&B_T^{-1}=A_T^{-1}\quad(\text{vector}),\\0,&B_T^{-1}=A_T^t\quad(\text{covector}).\end{cases}
\]
For a letter
outside the cut the Jacobian is the identity. The coefficient
estimate holds for each conjugate and hence for the cut bound.

Adjoining vector letters to a cut composes its substitution
with polynomial substitutions on those separate letters.
It can only increase the least order. Removing them composes
with inverse substitutions and loses at most the sum of their
degrees by homogeneity in each removed letter.  They give the local cutoff losses
when the binomial approximation is implemented by divided
derivatives and multiplication by monomials of the same degree.
\end{proof}
